\section{Quasi-closed-shell (non-aufbau) radicals}\label{sec:qcs}

\begin{reviewclaim}
``In molecular systems with unpaired electrons, the violation of the aufbau principle in electron
configuration is reflected in the conversion of the SOMO and HOMO orbital energies, yielding a so-called
quasi-closed-shell state (Figure 10). Its main fingerprint is the electron attachment or removal pattern:
upon ionization an electron is removed from the higher energy doubly occupied MO, leaving another
unpaired electron behind, while upon addition new electron enters the former LUMO; essentially, SOMO does
not participate in these processes.'' (p.~448--449) ``\dots a rigorous support for the non-aufbau
occupation can be obtained a posteriori from the ionization potentials and electron affinities \dots as
well as from the determination of the spin state \dots of the corresponding open-shell
oxidation/reduction products.'' (p.~449)
\end{reviewclaim}

\subsection{Frozen-orbital model of the redox fingerprint}
Three relevant orbitals: SOMO $s$, the highest doubly occupied orbital $h$, and the LUMO $l$. Ground
configuration $(\text{core})\,s^1h^2$ (a doublet).
\begin{enumerate}
  \item UHF orbital energies, from \cref{eq:koopmans} applied to each spin orbital (Slater--Condon, \cref{sec:koopmans}):
        \begin{align}
          \varepsilon_h^\alpha&=h_h+J_{hh}+J_{hs}-K_{hs}, &\varepsilon_h^\beta&=h_h+J_{hh}+J_{hs},\label{eq:uhfh}\\
          \varepsilon_s^\alpha&=h_s+2J_{hs}-K_{hs}, &\varepsilon_s^\beta(\text{virtual})&=h_s+J_{ss}+2J_{hs}-K_{hs}.\label{eq:uhfs}
        \end{align}
        (Interactions with the core are absorbed into $h_h,h_s$.)
  \item \textbf{Spin polarisation of the donor HOMO (Scheme~3).} From \cref{eq:uhfh},
        $\varepsilon_h^\beta-\varepsilon_h^\alpha=K_{hs}\ge0$: the $\alpha$ electron of $h$ (parallel to the unpaired spin) is
        stabilised by exchange with the SOMO electron, while the $\beta$ electron is not. The donor's $\alpha$ and
        $\beta$ HOMOs split. That is the content of Scheme~3.
  \item \textbf{Condition for SOMO--HOMO conversion} (within the $\alpha$ manifold):
        $\varepsilon_s^\alpha<\varepsilon_h^\alpha\iff h_s+2J_{hs}-K_{hs}<h_h+J_{hh}+J_{hs}-K_{hs}\iff h_s+J_{hs}<h_h+J_{hh}$.
        In words: the SOMO's one-electron energy plus its repulsion with \emph{one} $h$ electron lies below
        $h$'s one-electron energy plus its \emph{self}-repulsion. A stable, low-lying radical orbital and a
        separate, electron-rich (high $h_h$, e.g.\ anionic) fragment favour this.
  \item \textbf{Oxidation.} Koopmans: $\mathrm{IP}(h)=-\varepsilon_h<\mathrm{IP}(s)=-\varepsilon_s$, so the electron leaves $h$.
        The product $(\text{core})\,s^1h^1$ is a diradical cation (or diradical, if the parent was an anion). Its
        triplet and open-shell singlet are split by $E_{\text{OSS}}-E_T=2K_{sh}$ (\cref{eq:EOSS} with $a,b\to s,h$).
        With $K_{sh}>0$ the triplet is lower: the observed ``triplet cation diradicals with nondegenerate SOMOs''.
  \item \textbf{Reduction.} The added electron can enter the $\beta$ SOMO, at cost $\varepsilon_s^\beta$ of \cref{eq:uhfs}, giving
        a closed-shell anion; or the LUMO, at $\varepsilon_l$, giving the diradical $s^1h^2l^1$. Koopmans picks the
        lower: the diradical forms iff $\varepsilon_l<\varepsilon_s^\beta=h_s+J_{ss}+2J_{hs}-K_{hs}$. The large on-site
        repulsion $J_{ss}$ makes this possible. Note that this is a \emph{separate} condition, not implied by
        step 3; the review's Figure~10 depicts the case where it holds, and Figure~17 gives examples.
  \item \textbf{Protonation switch.} Protonating the anionic fragment strongly lowers $h_h$ (Coulomb attraction of
        $\mathrm H^+$). Once $h_h+J_{hh}<h_s+J_{hs}$ the aufbau order is restored, and oxidation removes the SOMO
        electron, giving a closed-shell cation. This is the pH-switchable conversion of Figure~16.
\end{enumerate}

\begin{caveatbox}
Steps 1--5 use frozen orbitals and canonical orbital energies. By \cref{sec:rohf} orbital energies are
convention-dependent, and by \cref{sec:koopmans} relaxation shifts them. As the review says, the
\emph{observable} criterion is energetic: compare the two possible oxidised states directly,
$E(\text{diradical},T)$ vs.\ $E(\text{closed shell})$ ($\Delta$SCF).
\end{caveatbox}

% numbers from verification/logs/t11.log (T11b) and logs/t11c.log (T11c)
\subsection{Numerical test: a distonic radical anion and its conjugate acid}
Model: the aminoxyl radical $^\bullet$ON(Me)CH$_2$CH$_2$COO$^-$, with a carboxylate that is not conjugated to
the radical, and its protonated form $^\bullet$ON(Me)CH$_2$CH$_2$COOH. UB3LYP/6-31+G*; radical geometries
optimised (geomeTRIC). Vertical oxidation at the radical geometry to (i) the triplet (one electron removed
from below the SOMO, the diradical) and (ii) the closed-shell singlet (SOMO electron removed).
\begin{center}
\begin{tabular}{lcccc}
\toprule
radical & $\alpha$-SOMO position & IE$\to$triplet & IE$\to$closed shell & $\Delta$(T$-$CS)\\
\midrule
acid (COOH) & is the $\alpha$-HOMO & 9.489 eV & 8.179 eV & $+1.31$ eV\\
anion (COO$^-$) & 3 levels below the $\alpha$-HOMO & 4.008 eV & 3.946 eV & $+0.06$ eV\\
\bottomrule
\end{tabular}
\end{center}
\begin{enumerate}
  \item \emph{Orbital level: confirmed.} In the anion the $\alpha$-SOMO ($-2.40$~eV) lies below three doubly
        occupied levels (character not analysed) ($\alpha$-HOMO $-1.62$~eV). In the acid the SOMO is the HOMO. This is
        the pH-switchable SOMO--HOMO conversion of step 6.
  \item \emph{Energetic level: shifted, but no crossover.} Deprotonation moves the
        triplet--closed-shell preference by $1.25$~eV towards the diradical, but at this level the closed-shell
        product of the anion remains lower, by $0.06$~eV. For this short-linker model, the orbital inversion
        alone does not guarantee that oxidation yields the diradical. This is exactly why the review insists on
        the energetic, observable criterion (caveat above). T11c below shows that this verdict depends on the
        functional.
\end{enumerate}
% T11c results (functional and separation dependence)
\emph{(T11c pending.)}



\begin{verifybox}
T2: \cref{eq:uhfh,eq:uhfs}, $\varepsilon_h^\beta-\varepsilon_h^\alpha=K_{hs}$, and the cation $E_{\text{OSS}}-E_T=2K_{hs}$ verified
numerically against Slater--Condon determinant energies for random integrals. \textbf{PASS}.
\end{verifybox}
\begin{sourcebox}
Gryn'ova, Marshall, Blanksby \& Coote, \emph{Nat.\ Chem.}\ \textbf{5}, 474 (2013); Gryn'ova \& Coote,
\emph{J.\ Am.\ Chem.\ Soc.}\ \textbf{135}, 15392 (2013): summaries only (SOMO--HOMO conversion when a stabilised
radical is not $\pi$-conjugated with an anionic group; protonation restores normal ordering). Westcott et al.\
(2000): VO(oep), where the unpaired metal d electron is bound more strongly than several ligand pairs
(author summary).
\end{sourcebox}
